% ═══════════════════════════════════════════════════════════════
% LERNBLATT: Formeln praktisch anwenden
% Informatik Klasse 7 | Stunde 2
% ═══════════════════════════════════════════════════════════════

\documentclass[11pt, a4paper]{article}

\usepackage[utf8]{inputenc}
\usepackage[T1]{fontenc}
\usepackage[ngerman]{babel}

\usepackage{helvet}
\renewcommand{\familydefault}{\sfdefault}

\usepackage[margin=1.5cm, top=1.8cm, bottom=1.5cm]{geometry}
\usepackage{enumitem}
\usepackage{tabularx}
\usepackage{booktabs}
\usepackage{array}
\usepackage{tikz}
\usepackage{fancyhdr}
\usepackage{xcolor}
\usepackage{tcolorbox}
\tcbuselibrary{skins}

% ═══════════════════════════════════════════════════════════════
% FARBEN
% ═══════════════════════════════════════════════════════════════

\definecolor{informatik}{HTML}{b35c00}
\definecolor{hellgrau}{HTML}{F5F5F5}
\definecolor{mittelgrau}{HTML}{888888}
\definecolor{excelgruen}{HTML}{217346}

\colorlet{fachfarbe}{informatik}

% ═══════════════════════════════════════════════════════════════
% KOPFZEILE
% ═══════════════════════════════════════════════════════════════

\pagestyle{fancy}
\fancyhf{}
\renewcommand{\headrulewidth}{0.4pt}
\renewcommand{\footrulewidth}{0pt}
\lhead{\small Informatik Klasse 7 | Tabellenkalkulation}
\rhead{\small Name: \underline{\hspace{3cm}}}

% ═══════════════════════════════════════════════════════════════
% BOXEN
% ═══════════════════════════════════════════════════════════════

\newtcolorbox{merkkasten}[1][]{
    colback=hellgrau,
    colframe=hellgrau,
    boxrule=0pt,
    arc=0pt,
    left=8pt, right=8pt, top=8pt, bottom=6pt,
    borderline north={2pt}{0pt}{fachfarbe},
    before skip=8pt, after skip=8pt,
    #1
}

\newtcolorbox{formelkasten}{
    colback=white,
    colframe=fachfarbe,
    boxrule=1pt,
    arc=0pt,
    left=10pt, right=10pt, top=8pt, bottom=8pt,
    before skip=6pt, after skip=6pt
}

\newtcolorbox{excelbox}{
    colback=white,
    colframe=excelgruen,
    boxrule=1pt,
    arc=0pt,
    left=8pt, right=8pt, top=6pt, bottom=6pt,
    before skip=6pt, after skip=6pt
}

\setlength{\parindent}{0pt}
\setlength{\parskip}{4pt}

% ═══════════════════════════════════════════════════════════════
% DOKUMENT
% ═══════════════════════════════════════════════════════════════

\begin{document}

% Header
\begin{center}
    {\Large\bfseries\textcolor{fachfarbe}{Formeln in der Tabellenkalkulation}}\\[0.3em]
    {\small Lernblatt | Stunde 2}
\end{center}
\vspace{-0.3em}
\hrule
\vspace{0.8em}

% ═══════════════════════════════════════════════════════════════
% 1. FORMELN
% ═══════════════════════════════════════════════════════════════

\textbf{\textcolor{fachfarbe}{1. Formeln eingeben}}

\begin{merkkasten}
\textbf{Merke:} Jede Formel beginnt mit dem \textbf{Gleichheitszeichen} \texttt{=}
\end{merkkasten}

\textbf{Grundrechenarten in Formeln:}

\begin{center}
\begin{tabular}{c c c l}
\toprule
\textbf{Mathematik} & \textbf{Taste} & \textbf{Beispiel} & \textbf{Ergebnis} \\
\midrule
$+$ (plus) & \texttt{+} & \texttt{=A1+B1} & Addiert A1 und B1 \\
$-$ (minus) & \texttt{-} & \texttt{=A1-B1} & Subtrahiert B1 von A1 \\
$\times$ (mal) & \texttt{*} & \texttt{=A1*B1} & Multipliziert A1 mit B1 \\
$\div$ (geteilt) & \texttt{/} & \texttt{=A1/B1} & Dividiert A1 durch B1 \\
\bottomrule
\end{tabular}
\end{center}

\vspace{0.3em}

\begin{excelbox}
\textbf{Beispiel:} In Zelle C1 steht die Formel \texttt{=A1+B1}

\vspace{0.3em}
\begin{center}
\begin{tikzpicture}[scale=0.9]
    % Header
    \fill[gray!20] (0,0) rectangle (1,0.6);
    \fill[gray!20] (1,0) rectangle (3,0.6);
    \fill[gray!20] (3,0) rectangle (5,0.6);
    \fill[gray!20] (5,0) rectangle (7,0.6);
    \draw (0,0) rectangle (1,0.6);
    \draw (1,0) rectangle (3,0.6);
    \draw (3,0) rectangle (5,0.6);
    \draw (5,0) rectangle (7,0.6);
    \node at (0.5,0.3) {\small};
    \node at (2,0.3) {\small\textbf{A}};
    \node at (4,0.3) {\small\textbf{B}};
    \node at (6,0.3) {\small\textbf{C}};

    % Row 1
    \fill[gray!20] (0,-0.6) rectangle (1,0);
    \draw (0,-0.6) rectangle (1,0);
    \draw (1,-0.6) rectangle (3,0);
    \draw (3,-0.6) rectangle (5,0);
    \draw[excelgruen, thick] (5,-0.6) rectangle (7,0);
    \node at (0.5,-0.3) {\small\textbf{1}};
    \node at (2,-0.3) {\small 10};
    \node at (4,-0.3) {\small 5};
    \node at (6,-0.3) {\small\textbf{15}};

    % Pfeil
    \draw[->, thick, fachfarbe] (7.3,-0.3) -- (8.5,-0.3);
    \node[right] at (8.5,-0.3) {\small\texttt{=A1+B1}};
\end{tikzpicture}
\end{center}
\end{excelbox}

\vspace{0.5em}

% ═══════════════════════════════════════════════════════════════
% 2. SUMME-FUNKTION
% ═══════════════════════════════════════════════════════════════

\textbf{\textcolor{fachfarbe}{2. Die SUMME-Funktion}}

\begin{formelkasten}
\centering
\large\texttt{=SUMME(Startadresse:Endadresse)}
\end{formelkasten}

\begin{itemize}[leftmargin=1.5em, itemsep=2pt]
    \item Addiert alle Werte in einem \textbf{Bereich}
    \item Der Doppelpunkt \texttt{:} bedeutet „von ... bis"
    \item Beispiel: \texttt{=SUMME(B2:B5)} addiert B2, B3, B4 und B5
\end{itemize}

\vspace{0.5em}

% ═══════════════════════════════════════════════════════════════
% 3. REFERENZPRINZIP
% ═══════════════════════════════════════════════════════════════

\textbf{\textcolor{fachfarbe}{3. Das Referenzprinzip}}

\begin{merkkasten}
\textbf{Referenzprinzip:} Ändert sich der Wert einer Zelle, werden alle Formeln, die diese Zelle verwenden, \textbf{automatisch neu berechnet}.
\end{merkkasten}

\begin{center}
\begin{tikzpicture}[scale=0.85]
    % Vorher
    \node[anchor=south] at (2,-0.8) {\small\textbf{Vorher:}};
    \fill[gray!20] (0,-1.4) rectangle (1,-0.8);
    \draw (0,-1.4) rectangle (1,-0.8);
    \draw (1,-1.4) rectangle (2.5,-0.8);
    \draw (2.5,-1.4) rectangle (4,-0.8);
    \node at (0.5,-1.1) {\small\textbf{1}};
    \node at (1.75,-1.1) {\small 10};
    \node at (3.25,-1.1) {\small 15};

    % Pfeil
    \draw[->, very thick, fachfarbe] (4.5,-1.1) -- (5.5,-1.1);
    \node[above] at (5,-1.1) {\tiny A1 ändern};

    % Nachher
    \node[anchor=south] at (7.5,-0.8) {\small\textbf{Nachher:}};
    \fill[gray!20] (5.8,-1.4) rectangle (6.8,-0.8);
    \draw (5.8,-1.4) rectangle (6.8,-0.8);
    \draw[fachfarbe, thick] (6.8,-1.4) rectangle (8.3,-0.8);
    \draw (8.3,-1.4) rectangle (9.8,-0.8);
    \node at (6.3,-1.1) {\small\textbf{1}};
    \node at (7.55,-1.1) {\small\textbf{20}};
    \node at (9.05,-1.1) {\small\textbf{25}};

    % Beschriftung
    \node[below, text width=4cm, align=center] at (2,-1.5) {\tiny C1 zeigt 15\\(= 10 + 5)};
    \node[below, text width=4cm, align=center] at (8,-1.5) {\tiny C1 zeigt automatisch 25\\(= 20 + 5)};
\end{tikzpicture}
\end{center}

\vspace{0.3em}

% ═══════════════════════════════════════════════════════════════
% 4. AUFZIEHEN (FILL-HANDLE)
% ═══════════════════════════════════════════════════════════════

\textbf{\textcolor{fachfarbe}{4. Der Aufzieh-Trick (Fill-Handle)}}

\begin{merkkasten}
\textbf{Profi-Tipp:} Formeln müssen nicht einzeln eingegeben werden! Mit dem \textbf{grünen Punkt} (Fill-Handle) kannst du Formeln automatisch kopieren.
\end{merkkasten}

\textbf{So funktioniert es:}
\begin{enumerate}[leftmargin=1.5em, itemsep=2pt]
    \item Formel in erste Zelle eingeben (z.B. \texttt{=B2*C2} in D2)
    \item Zelle anklicken $\rightarrow$ \textcolor{excelgruen}{\textbf{grüner Punkt}} erscheint rechts unten
    \item Punkt nach unten ziehen $\rightarrow$ Formel wird automatisch angepasst!
\end{enumerate}

\begin{excelbox}
\textbf{Beispiel:} Aus \texttt{=B2*C2} wird beim Aufziehen automatisch:\\[0.3em]
\hspace*{1em}\texttt{=B3*C3} $\rightarrow$ \texttt{=B4*C4} $\rightarrow$ \texttt{=B5*C5} $\rightarrow$ ...
\end{excelbox}

\vspace{0.3em}

% ═══════════════════════════════════════════════════════════════
% 5. WARUM IST DAS NÜTZLICH?
% ═══════════════════════════════════════════════════════════════

\textbf{\textcolor{fachfarbe}{5. Warum ist Tabellenkalkulation so mächtig?}}

\begin{itemize}[leftmargin=1.5em, itemsep=2pt]
    \item \textbf{Automatisierung:} Eine Formel kann 50.000 Berechnungen ersetzen
    \item \textbf{Echtzeit:} Ändere einen Wert $\rightarrow$ alles wird sofort neu berechnet
    \item \textbf{Keine Fehler:} Formeln vergessen nie zu rechnen
\end{itemize}

{\small\textcolor{mittelgrau}{Profis nutzen Excel für: Finanzen, Wissenschaft, Sportanalysen, Gaming-Stats, ...}}

% ═══════════════════════════════════════════════════════════════
% ZUSAMMENFASSUNG
% ═══════════════════════════════════════════════════════════════

\vspace{0.5em}

\begin{tcolorbox}[colback=hellgrau, colframe=fachfarbe, boxrule=1pt, arc=0pt, title={\textbf{Zusammenfassung}}, fonttitle=\bfseries\color{white}, coltitle=white, colbacktitle=fachfarbe]
\begin{itemize}[leftmargin=1.5em, itemsep=2pt]
    \item Formeln beginnen \textbf{immer} mit \texttt{=}
    \item Rechenzeichen: \texttt{+} \texttt{-} \texttt{*} \texttt{/}
    \item \texttt{=SUMME(B2:B5)} addiert einen Bereich
    \item Werte ändern $\rightarrow$ Formeln rechnen automatisch neu
    \item \textbf{Aufziehen:} Grüner Punkt kopiert Formeln automatisch
\end{itemize}
\end{tcolorbox}

\end{document}
